\documentclass[10pt,a4paper]{amsart}
\usepackage[margin=19mm]{geometry}
\usepackage{amsmath,amssymb,mathtools}
\usepackage[scr=rsfs,cal=euler]{mathalfa}
\usepackage{xfrac}
\usepackage[unicode,hidelinks,bookmarks=false]{hyperref}
\newcommand{\ie}{i.e.\ }
\newcommand{\cf}{cf.\ }
\newcommand{\resp}{respectively\ }
\newcommand{\Subsection}{\subsection}
\providecommand{\nobreakdash}{\nobreak\hbox{-}}
\setlength{\emergencystretch}{2em}
\multlinegap0pt
\usepackage{amssymb}
\usepackage{enumerate}
\usepackage{tikz}
\newtheorem{prop}{Proposition}
\newtheorem{claim}{Claim}
\newcommand{\N}{\mathbb N}
\newcommand\mcal{\mathcal}
\title{Complexity of Finite Borel Asymptotic Dimension}
\author{Jan Greb\'ik, Cecelia Higgins}
\date{Source: arXiv:2411.08797v3 (CC BY 4.0)}
\begin{document}
\maketitle

\noindent\emph{Source and licence.} Complexity of Finite Borel Asymptotic Dimension. Version arXiv:2411.08797v3 (CC BY 4.0), distributed by arXiv under the Creative Commons Attribution 4.0 International licence, http://creativecommons.org/licenses/by/4.0/, as recorded in the arXiv licence metadata for this exact version and shown on the abstract page https://arxiv.org/abs/2411.08797v3. Full licence notice: \texttt{licenses/arxiv-2411.08797-CC-BY-4.0.txt}.\par\smallskip

\noindent\textit{Editorial scope.} Counted unit: the article's prop pr:FIHtoCov (source0.tex lines 623--627). The article's complete original proof of it is delivered with it (source0.tex lines 628--674, one \texttt{proof} environment of 47 lines). No mathematics is written, repaired or re-derived: the statement and the proof are the article's own lines, in the article's own notation, and no retained line is substituted or re-flowed.  Retained uncounted as the source-local context the proof needs: lines 601--603.  Nothing was omitted from any retained span, so the package carries no omission record.  No retained line contains a citation, so the wrapper carries no bibliography and no bibliography entry.  No retained line contains a figure environment, an image-inclusion command or drawing code, so no image is delivered and the package declares no asset dependency.  Editorial formatting only, in addition: an AMS \texttt{amsart} preamble that restates the article's own declarations and macros used by the retained lines, namely the environments \texttt{prop}, \texttt{claim} on separate counters, together with the standard packages those lines need.  The retained environments print in this document as Proposition 1, Claim 1 in order of appearance, whereas the article prints the same items with the numbering of its own full document.  The complete native source of the article as distributed in the arXiv e-print for this exact version is bundled unchanged as \texttt{sources/168479/source0.tex}.
    \begin{equation}\label{eq:Decomp}
        \left\{0,1,\dots,\frac{r}{2}-1\right\}=I_0\sqcup \dots\sqcup I_{2m},
    \end{equation}

    \begin{prop}\label{pr:FIHtoCov}
        Let $t\in \N^+$, $f:X\to X$ be an acyclic Borel function and $H$ be a Borel $4(6t)^2$-forward independent hitting set.
        Define $c^{-1}_H(\{i\})=U_i$ for each $i\in \{0,1\}$, where $c_H$ is defined as in the previous paragraph.
        Then $U_0$ and $U_1$ are Borel sets such that $X=U_0\sqcup U_1$ and for each $i \in \{0, 1\}$, every equivalence class of $\mcal{F}_t(U_i)$ has diameter bounded by $28t+7$.
    \end{prop}

    \begin{proof}
        Recall from the definition of $c_H$ that $s=6t$.
        Define $\ell(x)\in \N^+$ to be minimal such that $c_H(x)\not=c_H(f^{\ell(x)}(x))$ and set
        $$e(x)=f^{\frac{s}{3}+\ell(x)}(x)$$
        for every $x\in X$.
        Note that we have $\ell(x)\le 2s+2$ for every $x\in X$.
        Indeed, if $k(x) > s + 1$, then by the definition of $c_{H}$, we have $\ell(x) \leq s + 1$.
        If $k(x) \leq s + 1$, then we have that $k(y) \ge r \geq s + 1$, where $y=f^{k(x)}(x)\in H$.        
        Hence, $\ell(f^{k(x)}(x)) \leq s + 1$ and consequently $\ell(x) \leq 2s + 2$.
        
        \begin{claim}\label{cl:BasicThree}
            For every $x\in U_i$, where $i\in \{0,1\}$, we have that
            $$\bigcup_{j=0}^{\frac{s}{3}} f^{-j}(e(x))\subseteq U_{1-i}.$$
        \end{claim}
        \begin{proof}
            As $e(x)=e(f^{\ell(x)-1}(x))$, we may assume without loss of generality that $\ell(x)=1$.
            First, we claim that $f^j(x)\in U_{1-i}\setminus H$ for every $1\le j\le \frac{s}{3}+1$.

            Indeed, if $f^j(x)\in H$ for any $1\le j\le \frac{s}{3}+1$, then, by the definition of $c_H$, we have $x\in U_{1-i}$, a contradiction.
            Similarly, assume for contradiction that $1\le j\le \frac{s}{3}+1$ is minimal such that $f^j(x)\not\in U_{1-i}$.
            Then $j\ge 2$, $f^{j-1}(x)\in U_{1-i}$ and we have that $x\in U_{1-i}$ as either $\lfloor k(x)/s \rfloor=\lfloor k(f^{j-1}(x))/s \rfloor \mod 2$ or $k(x),k(f^{j-1}(x))\in I_k$ for some $0\le k\le 2\ell$, where $I_k$ is from \eqref{eq:Decomp}.
            
            \medskip

            We finish the proof by showing that for every $0\le j\le \frac{s}{3}$ and $z\in f^{-j}(e(x))$, we have that $z\in U_{1-i}$.
            This clearly holds when $k(z)=k(f^{\frac{s}{3}+1-j}(x))$ by the previous paragraph.
            In particular, as $H$ is $r$-forward independent, it holds whenever $z\in H$.
            Now, the claim follows as, by the definition of $c_H$, we have that $c_H(w)=c_{H}(z)$ whenever $f^m(w)=z\in H$ and $0\le m\le s-1$.           
        \end{proof}
        
        \begin{claim}\label{cl:BasicTwo}
            Let $x,y\in X$ be such that $y\in [x]_{\mcal{F}_t(U_i)}$ for some $i\in\{0,1\}$.
            Then there is $k\in \N$ such that $f^k(y)=e(x)$.
        \end{claim}
        \begin{proof}
            Assume for contradiction that $f^k(y)\not= e(x)$ for any $k\in \N$.
            Then, by the definition of $\mcal{F}_t(U_i)$, there are $z,w\in [x]_{\mcal{F}_t(U_i)}$ such that $\rho_f(z,w)\le t$, $f^m(z)=e(x)$ for some $m\in \mathbb{N}$, and $f^k(w)\not= e(x)$ for any $k\in \N$.
            This is a contradiction, since the shortest path that connects $z$ and $w$ in $\mathcal{G}_f$ must contain $e(x)$ and go through
            $$\bigcup_{i=0}^{\frac{s}{3}} f^{-i}(e(x)),$$
            which is a subset of $U_{1 - i}$ by Claim~\ref{cl:BasicThree}.
            Hence, we have $\rho_f(w,z)\ge \frac{s}{3}>t$, a contradiction.         
        \end{proof}

        Observe that it follows from Claim~\ref{cl:BasicTwo} together with the assumption that $f$ is acyclic that, if $(x, y) \in \mcal{F}_r(U_i)$, then
        $$\rho_f(x,y)\le \rho_f(x,e(x))+\rho_f(y,e(y))\le \frac{2s}{3}+\ell(x)+\ell(y).$$
        This finishes the proof as we have that $\ell(x)\le 2s+2$ holds for every $x\in X$.
    \end{proof}

\end{document}
